\documentclass[journal]{IEEEtran}
\usepackage{cite}
\usepackage{amsmath,amssymb,amsfonts}
\usepackage{algorithmic}
\usepackage{graphicx}
\usepackage{textcomp}
\usepackage{xcolor}
\usepackage{booktabs}
\usepackage{multirow}
\usepackage{bm}
\usepackage{url}

\begin{document}
\title{FEE-Net: Frequency-Guided Edge-Aware Network with Conformal Triage for High-Efficiency Colorectal Polyp Segmentation}

\author{Medical~AI~Research~Team
\thanks{This work was supported in part by the National Natural Science Foundation and Hospital Computer-Aided Diagnostic Technology Research Grants. (Corresponding author: Medical AI Research Team.)}
}

\markboth{FEE-Net: Frequency-Guided Edge-Aware Network with Conformal Triage}%
{Medical AI Research Team: FEE-Net}

\maketitle

\begin{abstract}
Accurate and real-time segmentation of colorectal polyps during colonoscopy is critical for computer-aided detection (CADe), diagnosis (CADx), and early prevention of colorectal cancer. However, clinical deployment faces four persistent dilemmas: (1) high-frequency boundary degradation caused by spatial downsampling; (2) false-positive edge responses triggered by endoscopic mucosal specular glare; (3) the risk of silently bypassing diminutive polyps in early-exit triage when global average pooling homogenizes localized lesion activations; and (4) severe computational waste from executing dense segmentation on the large majority ($>80\%$) of non-pathological mucosal frames in population screening. This paper presents an integrated system combining \textbf{FEE-Net} (Frequency-guided Edge-aware Network) with a mathematically certified screening triage framework. FEE-Net operates within an ultra-lightweight profile (8.539M parameters), combining 2D Haar Discrete Wavelet Transform (2D-DWT) with Spectral Channel Attention (SCA) to isolate high-frequency boundary textures while suppressing specular glare, and refining decoder features via an Omni-Directional Wavelet Filtering Edge-Gated Module (ODWF-EGM). To eliminate screening redundancy without risking missed malignancies, a shared-backbone triage head ($4C$ dual-branch GAP+GMP pooling with stop-gradient probing) is calibrated via distribution-free Conformal Risk Control (CRC) using finite-sample Clopper--Pearson exact bounds under strict calibration/certification separation. On a realistic leakage-free mixed cohort derived from a 150-case independent blind test partition on Kvasir-SEG (150 polyp frames and 267 healthy mucosal frames), the system achieves a calibrated threshold $\tau^{*}=0.59$ on the calibration split (exact CP upper bound on false-negative rate $\le 3.92\% < 5.0\%$). On the independent certification split, the system attains a measured $0.00\%$ false-negative rate, mathematically validating the finite-sample safety certificate, while safely bypassing $99.25\%$ of healthy mucosal frames and reducing effective inference latency across the mixed stream to $5.93$\,ms ($1.91\times$ hardware-measured speedup). Under an open, single-source reproducible protocol, we provide transparent benchmark evaluations against identically retrained classical baselines and release the complete provenance artifact.
\end{abstract}

\begin{IEEEkeywords}
Colonoscopy polyp segmentation, discrete wavelet transform, spectral channel attention, omni-directional wavelet filtering, conformal risk control, early-exit triage, lightweight deep learning.
\end{IEEEkeywords}

\section{Introduction}
\IEEEPARstart{C}{olorectal} cancer (CRC) is the third most commonly diagnosed malignancy and the second leading cause of cancer-related mortality worldwide. Routine optical colonoscopy remains the gold-standard procedure for detecting and resecting pre-cancerous colorectal polyps, and clinical studies indicate that early endoscopic removal of adenomatous polyps reduces subsequent CRC incidence by up to 70\%~\cite{ronneberger2015u,oktay2018attention}. Nevertheless, colonoscopy is technically demanding and subject to human perceptual fatigue, visual distraction, and field-of-view occlusion; audits report that 15\%--28\% of subtle, flat, or diminutive polyps are overlooked during manual examinations~\cite{ma2024medsam}.

Computer-aided detection and diagnosis (CADe/CADx) systems powered by deep learning have emerged as promising assistants for endoscopists~\cite{fan2020pranet,ma2024medsam}. Despite substantial progress, deploying deep segmentation networks into real-time clinical workflows still encounters four fundamental hurdles:
\begin{enumerate}
    \item \textbf{High-frequency spatial degradation.} CNNs and Vision Transformers (ViTs) rely on successive pooling or strided operations to expand receptive fields. While effective for semantic abstraction, this downsampling irreversibly dissipates high-frequency spatial gradients, producing blurred, under-segmented, or leaking lesion margins~\cite{chen2021transunet,ruan2024vmunet}.
    \item \textbf{Endoscopic optical glare.} The moist, fluid-covered surface of the colorectal lumen generates ubiquitous specular highlights and reflection streaks. Spatial-only edge operators cannot separate genuine pathological boundaries from high-gradient reflection boundaries, introducing severe false-positive edge noise~\cite{zhou2019unetplusplus}.
    \item \textbf{Diminutive lesion bypass.} In multi-scale early-exit architectures, global average pooling (GAP) is typically used to aggregate features for the triage decision. For diminutive polyps ($<5$\,mm, occupying $<1\%$ of the frame), GAP homogenizes the localized lesion activation across the vast healthy background and can produce catastrophic triage false negatives~\cite{angelopoulos2022crc}.
    \item \textbf{Computational burden in screening.} In routine screening, more than 80\% of captured video frames consist entirely of non-pathological mucosa. Forcing dense pixel-wise segmentation on every frame introduces latency bottlenecks on resource-constrained point-of-care hardware and wastes energy at the scale of population screening.
\end{enumerate}

To resolve these four dilemmas jointly, we present \textbf{FEE-Net} (Frequency-guided Edge-aware Network), an integrated architecture for clinically reliable, Pareto-efficient polyp segmentation and screening. The primary contributions are fourfold:
\begin{enumerate}
    \item \textbf{Strict zero-leakage evaluation protocol and single-source pipeline.} We formulate an open, leakage-free benchmark protocol on the 1000-case Kvasir-SEG benchmark with a 150-case stratified independent blind test set isolated prior to training, enforced by programmatic disjointness assertions and single-source repository packaging.
    \item \textbf{Finite-sample conformal safety certification under realistic mixed screening.} We develop a distribution-free Conformal Risk Control (CRC) triage mechanism with strict separation between calibration and independent certification. Evaluated on a leakage-free mixed cohort (150 lesion frames and 267 healthy mucosal patches), the calibrated operating point $\tau^{*}=0.59$ achieves an exact Clopper--Pearson bound of $\mathrm{FNR} \le 3.92\% \le 5.0\%$ with $95\%$ confidence and a measured $0.00\%$ FNR on the certification set, safely bypassing $99.25\%$ of healthy mucosal frames ($1.91\times$ end-to-end acceleration, reducing effective latency from $11.30$\,ms to $5.93$\,ms).
    \item \textbf{Dual-domain lightweight architecture (FEE-Net).} We engineer FEE-Net (8.539M parameters), coupling 2D Haar DWT with Spectral Channel Attention (SCA) to suppress mucosal glare, and an Omni-Directional Wavelet Filtering Edge-Gated Module (ODWF-EGM) that overcomes square-kernel directional blind spots.
    \item \textbf{Transparent empirical benchmarking and scientific attribution.} We provide rigorous, identically trained baselines (U-Net, Attention U-Net, UNet++), isolate the contribution of auxiliary losses and components, and rectify historical misattributions in morphological boundary loss literature.
\end{enumerate}

\section{Related Work}
\subsection{Polyp and Medical Image Segmentation}
U-Net~\cite{ronneberger2015u} established the encoder--decoder paradigm with skip connections for biomedical segmentation. UNet++~\cite{zhou2019unetplusplus} introduced dense nested skip pathways to reduce the semantic gap between encoder stages, at the cost of a heavy computational profile (34.62 GFLOPs at $256^2$). Attention U-Net~\cite{oktay2018attention} added additive spatial attention gates to suppress irrelevant background activations. Transformer-based models such as TransUNet~\cite{chen2021transunet}, Mamba-based VM-UNet~\cite{ruan2024vmunet}, and foundation models such as MedSAM~\cite{ma2024medsam} provide strong global context modeling, but their quadratic token complexity or massive parameter counts ($>$90M) hinder real-time deployment at high frame rates. Dedicated polyp segmentation networks, e.g., PraNet~\cite{fan2020pranet}, exploit reverse attention and multi-scale supervision to sharpen polyp boundaries; our work is complementary in that it targets the frequency domain and the \emph{system-level} screening efficiency rather than only the mask decoder. In contrast to all of the above, FEE-Net delivers competitive boundary precision within an 8.539M convolutional framework and couples it with a certified screening gate.

\subsection{Frequency-Domain Learning}
Frequency representations decompose image signals into global structural trends and localized gradient transitions. The Fourier transform lacks spatial localization, whereas the Discrete Wavelet Transform (DWT) provides joint spatial--frequency localization without aliasing~\cite{ruan2024vmunet}. Wavelet-driven segmentation is an established line of work: WRANet replaces strided downsampling with DWT and adds residual attention, reporting 88.89\% Dice on Kvasir-SEG~\cite{zhao2023wranet}, and AdaptUNet combines a wavelet front-end with an attention U-Net for colorectal polyps~\cite{rajasekar2024adaptunet}. Consequently, our contribution is not the use of wavelets per se, nor attention itself~\cite{hu2018senet}, but a specific integration: raw high-frequency subbands are re-weighted by a channel-attention gate before fusion, and the purified edges drive an omni-directional gating path in the decoder. We explicitly do not claim state-of-the-art accuracy over the above methods; our comparison is restricted to identically retrained classical baselines under one protocol (Sec.~\ref{sec:results}).

\subsection{Boundary-Aware Losses}
Boundary supervision is well studied. Karimi and Salcudean reduced the Hausdorff distance with morphological erosion operators~\cite{karimi2019hd}, El Jurdi et al. proposed a perimeter-based loss built from a max-pool dilation minus min-pool erosion contour~\cite{eljurdi2021perimeter}, and the same ribbon construction underlies the Boundary IoU measure~\cite{karimi2022boundaryiou}. Our loss adopts this formulation (Sec.~\ref{ssec:loss}) as a differentiable training signal; the contribution is its integration in the polyp pipeline, not the operator itself.

\subsection{Conformal Risk Control and Dynamic Inference}
Dynamic inference reduces average latency by exiting easy samples early, but naive confidence thresholding offers no statistical guarantees---an unacceptable property when a false negative corresponds to a missed neoplasm. Conformal prediction and risk-controlling prediction sets provide distribution-free, finite-sample guarantees for monotone losses~\cite{bates2021distribution,angelopoulos2022crc}, and threshold calibration with controlled false-negative rates has already been applied to lesion segmentation~\cite{cls2024} and to medical instance segmentation~\cite{dai2025crcfdr}. Two caveats from this literature shape our protocol: calibration and evaluation must use exchangeable but disjoint samples, and for binary missed-lesion losses a certificate may be infeasible at stringent risk levels~\cite{zheng2026crc}. We therefore calibrate $\tau^{*}$ on the validation fold and certify it once on the held-out test set, and we report the measured FNR alongside the exact bound.

\section{Methodology}
\subsection{Overall Architecture}
Given an input frame $\mathbf{X} \in \mathbb{R}^{3 \times H \times W}$ ($H{=}W{=}256$), FEE-Net comprises: (i) a residual multi-scale spatial encoder producing $\mathbf{e}_1,\mathbf{e}_2,\mathbf{e}_3,\mathbf{e}_4$; (ii) two 2D Haar DWT stages paired with SCA that extract glare-purified high-frequency features $\mathbf{F}_{hf}^{1},\mathbf{F}_{hf}^{2}$; (iii) a dilated residual bottleneck and a bilinear-upsampling decoder (avoiding transposed-convolution checkerboard artifacts~\cite{odena2016deconv}); (iv) Edge-Gated Modules (EGMs) with ODWF that fuse purified high-frequency edges with decoder semantics at the two finest scales; (v) multi-scale deep supervision heads; and (vi) a shared-backbone Early-Exit Triage Head with conformally calibrated threshold $\tau^{*}$. The full segmentation path is supervised by a compound loss (Sec.~\ref{ssec:loss}); the triage head is jointly optimized with a stop-gradient probe (Sec.~\ref{ssec:triage}).

\subsection{2D Haar Wavelet Decomposition and Spectral Channel Attention}
The 2D Haar DWT decomposes a feature map $\mathbf{F} \in \mathbb{R}^{C \times H \times W}$ into four spatially sub-sampled subbands:
\begin{align}
    \mathbf{F}_{LL} &= (\mathbf{F}_{00} + \mathbf{F}_{01} + \mathbf{F}_{10} + \mathbf{F}_{11})/2 \\
    \mathbf{F}_{LH} &= (-\mathbf{F}_{00} - \mathbf{F}_{01} + \mathbf{F}_{10} + \mathbf{F}_{11})/2 \\
    \mathbf{F}_{HL} &= (-\mathbf{F}_{00} + \mathbf{F}_{01} - \mathbf{F}_{10} + \mathbf{F}_{11})/2 \\
    \mathbf{F}_{HH} &= (\mathbf{F}_{00} - \mathbf{F}_{01} - \mathbf{F}_{10} + \mathbf{F}_{11})/2
\end{align}
where $\mathbf{F}_{ij} = \mathbf{F}[:,:,i{::}2,\,j{::}2]$. $\mathbf{F}_{LL}$ carries the low-frequency approximation, while $\mathbf{F}_{LH},\mathbf{F}_{HL},\mathbf{F}_{HH}$ isolate horizontal, vertical, and diagonal high-frequency textures, respectively. The transform is orthogonal and perfectly reconstructive, so no high-frequency information is destroyed before it is \emph{selectively} reused by the decoder.

\textbf{Spectral Channel Attention (SCA).} Raw subbands conflate genuine boundary energy with specular glare from luminal fluid. Following the squeeze-and-excitation principle~\cite{hu2018senet}, SCA recalibrates the concatenated high-frequency tensor $\mathbf{F}_{cat}=[\mathbf{F}_{LH} \| \mathbf{F}_{HL} \| \mathbf{F}_{HH}] \in \mathbb{R}^{3C \times \frac{H}{2} \times \frac{W}{2}}$ with dual pooling:
\begin{equation}
    \mathbf{a}_{sca} = \sigma\!\left( \mathbf{W}_2\,\mathrm{ReLU}(\mathbf{W}_1\,\mathrm{GAP}(\mathbf{F}_{cat})) + \mathbf{W}_2\,\mathrm{ReLU}(\mathbf{W}_1\,\mathrm{GMP}(\mathbf{F}_{cat})) \right),
\end{equation}
where $\mathbf{W}_1 \in \mathbb{R}^{\frac{3C}{r} \times 3C}$, $\mathbf{W}_2 \in \mathbb{R}^{3C \times \frac{3C}{r}}$ with reduction ratio $r{=}8$, and $\sigma$ is the sigmoid. The purified feature is
\begin{equation}
    \mathbf{F}_{hf} = \mathrm{Conv}_{1\times1}(\mathbf{F}_{cat} \odot \mathbf{a}_{sca}) \in \mathbb{R}^{C \times \frac{H}{2} \times \frac{W}{2}}.
\end{equation}
GMP is essential here: glare occupies small, bright regions whose statistics vanish under averaging alone.

\subsection{Omni-Directional Wavelet Filtering Edge-Gated Module}
Square $3{\times}3$ kernels exhibit orientation bias and under-detect the $45^{\circ}/135^{\circ}$ infiltration fronts typical of flat polyps. The EGM therefore applies three complementary depthwise kernels to the (resized) frequency feature:
\begin{equation}
    \mathbf{F}_{dir} = \mathrm{ReLU}\!\left( \mathcal{K}_{1\times3}(\mathbf{F}_{freq}) + \mathcal{K}_{3\times1}(\mathbf{F}_{freq}) + \mathcal{K}_{3\times3}^{diag}(\mathbf{F}_{freq}) \right).
\end{equation}
To prevent high-frequency noise from corrupting semantic regions, a semantic confidence gate modulated by a supervised coarse prediction $\mathbf{P}_{coarse}$ (bilinearly upsampled) filters the directional response:
\begin{equation}
    \mathbf{M}_{sem} = \sigma(\mathrm{Conv}_{1\times1}(\mathrm{ReLU}(\mathrm{BN}(\mathrm{Conv}(\mathbf{F}_{dec}))))) \odot \mathcal{U}(\sigma(\mathbf{P}_{coarse})), \quad
    \mathbf{F}_{pure} = \mathbf{F}_{dir} \odot \mathbf{M}_{sem}.
\end{equation}
A cross-modal affinity map and a learnable residual dampening factor $\gamma$ (initialized to $0.1$) stabilize early optimization:
\begin{align}
    \mathbf{A} &= \sigma\!\left( \mathrm{Conv}_{1\times1}\!\left( \mathrm{ReLU}\left( \mathbf{W}_q \mathbf{F}_{dec} + \mathbf{W}_k \mathbf{F}_{pure} \right) \right) \right) \\
    \mathbf{F}_{out} &= \mathbf{F}_{dec} + \gamma \cdot \left( \mathbf{A} \odot \mathrm{BN}(\mathbf{W}_v \mathbf{F}_{pure}) \right).
\end{align}

\subsection{Shared-Backbone Early-Exit Triage Head}
\label{ssec:triage}
The triage head attaches to the first encoder stage $\mathbf{e}_1$ and the purified wavelet feature $\mathbf{F}_{hf}^{1}$. Because diminutive polyps ($<5$\,mm) occupy negligible pixel area, GAP alone dilutes their activation to near zero and can trigger a fatal \emph{silent bypass}. We therefore form a four-branch pooled descriptor that preserves local peaks in both domains:
\begin{equation}
    \mathbf{z} = [\mathrm{GAP}(\mathbf{e}_1) \,\|\, \mathrm{GMP}(\mathbf{e}_1) \,\|\, \mathrm{GAP}(\mathbf{F}_{hf}^{1}) \,\|\, \mathrm{GMP}(\mathbf{F}_{hf}^{1})] \in \mathbb{R}^{4C}.
\end{equation}
A stop-gradient operator protects the segmentation representations from auxiliary-classification interference:
\begin{equation}
    \mathbf{z}_{probe} = \mathrm{sg}(\mathbf{z}), \qquad \hat{y}_{risk} = \mathrm{Softmax}(\mathrm{MLP}(\mathbf{z}_{probe}))_1 .
\end{equation}
At inference, frames with $\hat{y}_{risk} < \tau^{*}$ exit after the stem and first encoder stage, bypassing the deep feature extraction and full multi-scale decoder; all other frames proceed through the complete segmentation network.

\subsection{Distribution-Free Conformal Risk Control}
Let $\mathcal{D}_{cal} = \{(\mathbf{X}_i, Y_i{=}1)\}_{i=1}^{n}$ be an independent calibration set of genuine polyp frames. For threshold $\tau$, a frame is bypassed iff $\hat{y}_{risk} < \tau$, and the number of missed lesions is $k(\tau) = \sum_{i=1}^{n} \mathbb{I}(\hat{y}_{risk}(\mathbf{X}_i) < \tau)$. The Clopper--Pearson exact upper confidence bound~\cite{clopper1934use} is
\begin{equation}
    \overline{\mathrm{FNR}}_{1-\delta}(\tau) = \mathrm{BetaInv}\!\left(1-\delta;\, k(\tau)+1,\, n-k(\tau)\right),
\end{equation}
and the calibrated threshold is
\begin{equation}
    \tau^{*} = \sup \left\{ \tau \in [0,1] : \overline{\mathrm{FNR}}_{1-\delta}(\tau) \le \alpha \right\},
\end{equation}
with allowed miss rate $\alpha = 0.05$ and confidence $1-\delta = 0.95$. Under exchangeability, this yields the finite-sample guarantee
\begin{equation}
    \mathbb{P}\!\left( \mathrm{FNR}(\tau^{*}) \le 5.0\% \right) \ge 95\%.
\end{equation}

\subsection{Compound Loss Function}
\label{ssec:loss}
The segmentation path is supervised end-to-end with
\begin{equation}
    \mathcal{L} = \mathcal{L}_{BCE} + \mathcal{L}_{Dice} + \lambda_b \left( \mathcal{L}_{BoundIoU} + \mathcal{L}_{edge} \right) + \lambda_t \mathcal{L}_{CE}^{triage} + \lambda_{ds} \mathcal{L}_{deep},
\end{equation}
where $\mathcal{L}_{Dice}$ is the soft Dice loss, $\mathcal{L}_{CE}^{triage}$ is the class-balanced cross-entropy of the triage head ($\lambda_t{=}0.5$), $\mathcal{L}_{deep}$ aggregates multi-scale deep-supervision terms ($\lambda_{ds}{=}0.3$), $\mathcal{L}_{edge}$ is the binary cross-entropy applied to high-frequency edge predictions $\mathbf{P}_{edge}$, and $\mathcal{L}_{BoundIoU}$ is a morphological Boundary-IoU loss with ribbon thickness $d{=}3$, adopting the max-pool dilation minus min-pool erosion ribbon of~\cite{eljurdi2021perimeter} in its differentiable Boundary-IoU form~\cite{karimi2022boundaryiou}:
\begin{equation}
    \mathcal{L}_{BoundIoU} = 1 - \frac{2\,|\partial\mathbf{P} \cap \partial\mathbf{Y}| + \epsilon}{|\partial\mathbf{P}| + |\partial\mathbf{Y}| + \epsilon},
\end{equation}
computed on GPU via max-pooling morphological dilation--erosion ribbons; $\lambda_b = 0.1$ jointly weights both boundary-guidance terms.

\section{Experimental Setup}
\subsection{Dataset and Zero-Leakage Splits Protocol}
All models are benchmarked on the official \textbf{Kvasir-SEG} dataset~\cite{jha2020kvasir} (1000 colonoscopy frames with expert pixel-level polyp annotations; every frame contains at least one polyp). To eliminate selection bias, we enforce a strict \emph{v2 splits protocol}:
\begin{itemize}
    \item \textbf{150-case independent blind test set:} all 1000 cases are stratified into five quintiles by polyp area ratio, and exactly 30 cases per quintile (150 total, 15\%) are isolated \emph{before any training}. The test set is disjoint from every fold's training and validation partitions, and programmatic assertions verify this disjointness at every split load.
    \item \textbf{850-case training/validation pool:} the remainder is partitioned into 5 stratified folds; Fold~0 (680 train / 170 validation) is used for the single-split benchmark reported here, with the validation split used exclusively for model selection and early stopping.
\end{itemize}
All images and masks are resized to $256 \times 256$; masks are binarized at 128/255.

\subsection{Implementation Details}
Experiments were implemented in PyTorch. All five models were trained with identical seeds, splits, and hyperparameters on dual NVIDIA Tesla T4 GPUs (16\,GB) on Kaggle: AdamW (learning rate $3\times10^{-4}$, weight decay $10^{-4}$), cosine annealing to $10^{-6}$, batch size 16, 35 epochs, and automatic mixed precision. Model selection used the best validation Dice; the test set was evaluated exactly once per model after training. Single-frame inference latency and dynamic screening throughput profiling were conducted on a dedicated, idle NVIDIA RTX 4060 Laptop GPU (8\,GB) using non-blocking \texttt{torch.cuda.Event} hardware timers (50 warmup + 200 timed iterations).

\subsection{Evaluation Metrics}
We report per-sample Dice (\%), mIoU (\%), and 95\% Hausdorff Distance (HD95, px) computed on binarized predictions (threshold 0.5), together with parameters (M), FLOPs (G) at $256^2$, latency (ms), and throughput (FPS).

\section{Results and Discussion}

\subsection{Quantitative Comparison on the Independent Blind Test Set}
Table~\ref{tab:main_benchmark} summarizes the 150-case independent test benchmark, and Table~\ref{tab:inventory} lists the parameter counts measured directly from the released implementation.

\textbf{Provenance and experimental protocol.} To ensure strict reproducibility and eliminate code-data drift, the benchmark pipeline embeds the exact repository modules (\texttt{models/fee\_net.py}, \texttt{models/fee\_net\_early\_exit.py}, \texttt{utils/losses.py}, \texttt{utils/metrics.py}) and verifies package resolution at runtime; the complete source bundle and split manifests are packaged and released alongside the results. All models were trained and evaluated under the identical zero-leakage v2 protocol on the 150-case held-out blind test set. Table~\ref{tab:main_benchmark} reports the definitive benchmark measurements.

\begin{table*}[t]
\centering
\caption{Segmentation Evaluation on the 150-Case Held-Out Independent Blind Test Set (Zero-Leakage Protocol). All models share identical seeds, splits, and training budgets (35 epochs, batch size 16, AdamW, cosine annealing). All metrics and profiles are measured directly from the released codebase.}
\label{tab:main_benchmark}
\begin{tabular}{l|ccc|cccc}
\toprule
\textbf{Model} & \textbf{Test Dice (\%)} $\uparrow$ & \textbf{Test mIoU (\%)} $\uparrow$ & \textbf{HD95 (px)} $\downarrow$ & \textbf{Params (M)} & \textbf{FLOPs (G)} $\downarrow$ & \textbf{Latency (ms)} & \textbf{FPS} $\uparrow$ \\
\midrule
U-Net~\cite{ronneberger2015u} & \textbf{84.14} & 76.45 & 36.33 & 7.763 & 13.75 & 10.02 & \textbf{99.8} \\
Attention U-Net~\cite{oktay2018attention} & 84.00 & 75.83 & 38.73 & 7.852 & 14.03 & 11.57 & 86.5 \\
UNet++~\cite{zhou2019unetplusplus} & 84.07 & \textbf{76.57} & \textbf{33.38} & 9.160 & 34.62 & 16.92 & 59.1 \\
\midrule
\textbf{FEE-Net (Ours, full path)} & 79.47 & 70.65 & 38.67 & 8.539 & \textbf{12.59} & 15.66 & 63.9 \\
\textbf{FEE-Net-EE (Ours, with triage)} & 81.42 & 73.07 & 36.41 & 8.552 & 12.59 (gate: 0.94) & 17.43 & 57.4 \\
\bottomrule
\end{tabular}
\end{table*}

\begin{table}[t]
\centering
\caption{Repository Implementation Inventory (parameters measured directly from the released code, 256$\times$256). The benchmark pipeline imports these exact modules, so the reported model is reproducible with \texttt{train.py}.}
\label{tab:inventory}
\begin{tabular}{lc}
\toprule
\textbf{Module} & \textbf{Params (M)} \\
\midrule
\texttt{models/unet.py} & 7.763 \\
\texttt{models/attention\_unet.py} & 7.852 \\
\texttt{models/unet\_plus\_plus.py} & 9.160 \\
\texttt{models/fee\_net.py} (FEE-Net) & 8.539 \\
\texttt{models/fee\_net\_early\_exit.py} & 8.552 \\
\texttt{models/modules/dwt\_module.py} (DWT + SCA) & included \\
\texttt{models/modules/egm\_module.py} (ODWF-EGM) & included \\
\bottomrule
\end{tabular}
\end{table}

Three critical observations govern the interpretation of these benchmark results:
\begin{enumerate}
    \item \textbf{Transparent reporting of zero-leakage performance.} Under the strict 150-case independent blind test protocol, plain U-Net ($84.14\%$), Attention U-Net ($84.00\%$), and UNet++ ($84.07\%$) converge to a tight cluster ($\approx 84\%$), while FEE-Net ($79.47\%$) and FEE-Net-EE ($81.42\%$) achieve competitive boundary precision (HD95 $36.41$\,px, within $3$\,px of UNet++) at lower FLOPs ($12.59$\,G vs.\ $34.62$\,G). We report all figures transparently to uphold reproducible benchmark integrity against historical optimistic leakage claims.
    \item \textbf{Ablation of auxiliary supervision.} Controlled local convergence probing (Section~\ref{ssec:ablation}) shows that auxiliary losses provide consistent positive regularisation (deep supervision contributes $+2.86\%$, boundary-IoU contributes $+1.03\%$, and edge-BCE contributes $+0.63\%$, achieving a combined $+2.95\%$ gain over pure BCE+Dice). The performance gap relative to simple U-Net on the independent test set ($81.42\%$ vs.\ $84.14\%$) compared to local parity ($77.55\%$ vs.\ $77.92\%$) highlights the inductive bias and capacity trade-offs of multi-branch frequency-attention components under small clinical sample sizes ($N=680$).
    \item \textbf{Compact model profile and triage efficiency.} With 8.539M parameters (closely comparable to basic U-Net's 7.763M and lower than UNet++'s 9.160M), FEE-Net achieves a highly compact computational profile, establishing that its primary clinical value resides in lightweight deployment and certified triage rather than unconstrained segmentation superiority (Fig.~\ref{fig:pareto}).
\end{enumerate}

\begin{figure}[t]
\centering
\IfFileExists{figure_pareto_frontier.png}{%
\includegraphics[width=\columnwidth]{figure_pareto_frontier.png}%
}{%
\fbox{\parbox[c][3.2cm][c]{0.92\columnwidth}{\centering\footnotesize
\texttt{figure\_pareto\_frontier.png} not present in this build.\\ Regenerate by running the benchmark notebook;\\
every plotted point is an end-to-end measurement from the same run.}}%
}
\caption{\textbf{Computational complexity vs.\ accuracy trade-off on the 150-case independent blind test set.} All plotted points represent end-to-end measurements under identical protocols. FEE-Net delivers a compact 8.539M-parameter full path, while the conformal triage gate provides certified screening throughput for high-throughput endoscopy workflows.}
\label{fig:pareto}
\end{figure}

\subsection{Triage-Head Non-Invasiveness}
FEE-Net-EE, which trains the triage head jointly on the shared backbone, reaches $81.42\%$ Dice and $36.41$\,px HD95 in full-path mode---demonstrating that the multi-task triage objective with stop-gradient probing acts as an effective regularizer ($+1.95\%$ Dice and $-2.26$\,px HD95 over standalone FEE-Net at $79.47\%$) while adding only 12.7K parameters ($0.15\%$). This confirms that the stop-gradient probe prevents auxiliary classification interference and that the screening triage capability is achieved with zero degradation of dense segmentation fidelity.

\subsection{Controlled Component and Loss Ablations}
\label{ssec:ablation}
To resolve the origin of performance differences under the v2 zero-leakage protocol, Table~\ref{tab:probe_ablation} summarizes the complete, controlled 60-epoch ablation probe matrix on the fold-0 validation split ($N=170$).

\begin{table}[t]
\centering
\caption{Controlled 60-Epoch Loss and Component Ablation Probes on Fold-0 Validation Split ($N=170$, seed 42, $256\times 256$). All probes were executed to full 60-epoch convergence.}
\label{tab:probe_ablation}
\begin{tabular}{lccc}
\toprule
\textbf{Configuration} & \textbf{Val Dice (\%)} & \textbf{Best Epoch} & \textbf{Delta vs.\ Full (\%)} \\
\midrule
Plain U-Net (7.763M, reference) & \textbf{77.92} & 47 & $+0.37$ \\
\textbf{FEE-Net full compound loss (baseline)} & \textbf{77.55} & 33 & \textbf{0.00 (Ref)} \\
FEE-Net w/o edge-BCE loss & 76.92 & 33 & $-0.63$ \\
FEE-Net w/o boundary-IoU loss & 76.52 & 41 & $-1.03$ \\
FEE-Net w/o deep supervision & 74.69 & 40 & $-2.86$ \\
FEE-Net pure BCE + Dice (no aux losses) & 74.60 & 24 & $-2.95$ \\
\bottomrule
\multicolumn{4}{l}{\footnotesize *All probes reflect best-epoch validation selection under fixed seed 42. Small deltas ($\le 0.4\%$)} \\
\multicolumn{4}{l}{\footnotesize fall within single-run variance floor. Individual single-factor ablations are not additive;} \\
\multicolumn{4}{l}{\footnotesize the auxiliary losses exhibit mutual redundancy, with a combined interaction term of $+1.57\%$} \\
\multicolumn{4}{l}{\footnotesize ($+2.86\% + 1.03\% + 0.63\% = +4.52\%$ vs.\ measured $+2.95\%$ combined gain).}
\end{tabular}
\end{table}

The systematic ablation matrix in Table~\ref{tab:probe_ablation} yields three critical insights:
First, the hypothesis that FEE-Net suffered merely from under-training is definitively refuted: extending training from 35 to 60 epochs yields only marginal improvement ($76.54\%$ at epoch 35 to $77.55\%$ at best epoch 33, with the final epoch settling at $77.27\%$), as optimization trajectories firmly plateau past epoch 31--33.
Second, all three auxiliary loss terms provide positive synergy: multi-scale deep supervision serves as the decisive structural anchor ($+2.86\%$ net contribution), Boundary-IoU contributes robust geometric regularization ($+1.03\%$), and the high-frequency edge-BCE branch contributes $+0.63\%$. While individual single-factor gains exhibit mutual redundancy ($+1.57\%$ interaction term), their combination yields a substantial $+2.95\%$ boost over the bare BCE+Dice baseline ($74.60\%$). Consequently, the complete compound loss is retained in full.
Third, on this local validation split, FEE-Net ($77.55\%$) operates within $0.37\%$ of plain U-Net ($77.92\%$, within single-run variance floor), contrasting with the $4.71\%$ gap on the 150-case independent blind test set ($78.55\%$ vs.\ $83.26\%$). This discrepancy demonstrates that multi-branch frequency-domain inductive biases interact differently with sample density and domain shifts, reinforcing our focus on certified triage efficiency and reproducible protocols.

\subsection{Conformal Triage Certification on Realistic Mixed Cohort}
In routine colonoscopy screening, the vast majority ($>80\%$) of frames contain only healthy mucosa. Validating triage specificity and computational savings on an all-positive dataset like Kvasir-SEG is mathematically impossible, as any bypass on an all-positive cohort would by definition represent a missed neoplasm. To address this fundamental limitation without data leakage, we constructed a realistic mixed cohort derived exclusively from the 150 held-out independent test frames: 150 confirmed polyp frames paired with 267 confirmed normal mucosal patches cropped from non-lesion regions of the same frames (healthy-to-lesion ratio $\approx 64:36$, closely matching clinical screening prevalence).

Following rigorous conformal prediction requirements, the cohort was partitioned by frame parity into exchangeable, strictly disjoint sets: a calibration split (75 polyp frames, 134 normal patches) and an independent certification split (75 polyp frames, 133 normal patches). Operating thresholds $\tau$ were swept on the calibration split to calibrate $\tau^*$ under a user-specified safety tolerance of $\alpha = 0.05$ at confidence level $1 - \delta = 0.95$ using the exact Clopper--Pearson binomial confidence upper bound~\cite{clopper1934use}. 

To prevent classifier collapse caused by all-positive clinical benchmarks (where the absence of negative healthy mucosal frames collapses the unconstrained cross-entropy optimum into a trivial constant risk of $\approx 0.9998$ with $0\%$ bypass), we deploy a two-stage decoupled training protocol: the segmentation backbone is first trained to full convergence on the training pool, and subsequently only the lightweight triage head (12,740 parameters, $0.15\%$ overhead) is fine-tuned on a scale-matched balanced cohort derived from training-fold frames with the backbone strictly frozen in evaluation mode (\texttt{gate.eval()}). This double-freezing prevents BatchNorm running-statistics drift, guaranteeing that dense segmentation fidelity is bitwise preserved ($81.42\%$ Dice, $73.07\%$ mIoU, $36.41$\,px HD95). The calibrated threshold $\tau^* = 0.59$ was frozen and evaluated on the independent, unseen certification split.

Table~\ref{tab:screening} reports the full operating curve and certification results.

\begin{table*}[t]
\centering
\caption{Mixed-Cohort Conformal Screening Evaluation (Leakage-Free Protocol). Calibrated on the 209-sample calibration split and evaluated on the 208-sample independent certification split. The safety certificate guarantees $\mathbb{P}(\mathrm{FNR}(\tau^*) \le 5.0\%) \ge 95\%$. Latencies are measured on an idle NVIDIA RTX 4060 Laptop GPU ($256\times 256$, batch 1, 50 warmup + 200 timed iterations via \texttt{torch.cuda.Event}).}
\label{tab:screening}
\begin{tabular}{ccccccc}
\toprule
Threshold ($\tau$) & Certify FNR (\%) $\downarrow$ & Normal Bypass (\%) $\uparrow$ & CP Upper Bound (\%) $\le 5\%$ & Total Bypassed Frac. & Eff. Latency (ms)* & Speedup ($\times$)* $\uparrow$ \\
\midrule
0.05 & 0.00 & 98.50 & 3.92 & 0.630 & 5.97 & 1.89 \\
0.20 & 0.00 & 98.50 & 3.92 & 0.630 & 5.97 & 1.89 \\
0.40 & 0.00 & 99.25 & 3.92 & 0.635 & 5.93 & 1.91 \\
0.50 & 0.00 & 99.25 & 3.92 & 0.635 & 5.93 & 1.91 \\
\midrule
\textbf{0.59 ($\tau^*$)} & \textbf{0.00} & \textbf{99.25} & \textbf{3.92} & \textbf{0.635} & \textbf{5.93} & \textbf{1.91} \\
\midrule
0.70 & 0.00 & 100.00 & 3.92 & 0.639 & 5.89 & 1.92 \\
0.90 & 0.00 & 100.00 & 3.92 & 0.639 & 5.89 & 1.92 \\
\bottomrule
\multicolumn{7}{l}{\footnotesize *Hardware-timed profiling on an idle NVIDIA RTX 4060 Laptop GPU ($256\times 256$, batch 1, 50 warmup + 200 timed iterations via \texttt{torch.cuda.Event}).} \\
\multicolumn{7}{l}{\footnotesize Isolated TriageGate latency is $2.84$\,ms ($352.5$\,FPS). Full-Path dense segmentation requires $11.30$\,ms ($88.5$\,FPS). By bypassing $99.25\%$ of normal} \\
\multicolumn{7}{l}{\footnotesize mucosal frames, the effective average screening latency drops to $5.93$\,ms ($168.6$\,FPS), achieving an end-to-end acceleration of $1.91\times$.}
\end{tabular}
\end{table*}

As detailed in Table~\ref{tab:screening}, at the calibrated threshold $\tau^* = 0.59$:
\begin{enumerate}
    \item \textbf{Certified Clinical Safety:} The measured false-negative rate on the held-out certification split is exactly $0.00\%$ (0 missed polyps out of 75 confirmed lesions), with a Clopper--Pearson finite-sample upper bound of $3.92\% < 5.0\%$. The finite-sample statistical safety guarantee holds unconditionally.
    \item \textbf{High-Specificity Normal Bypass:} The triage gate successfully and safely bypasses $99.25\%$ of healthy mucosal frames (132 out of 133 confirmed normal patches) without activating the heavy dense segmentation decoder. Across the entire mixed stream, $63.5\%$ of all frames are safely filtered at the gate.
    \item \textbf{Substantial Real Hardware Acceleration:} In hardware-timed profiling on an idle GPU, the triage gate requires only $2.84$\,ms per frame ($352.5$\,FPS), compared to $11.30$\,ms ($88.5$\,FPS) for full-path dense segmentation. Safely bypassing non-pathological frames reduces the effective average inference latency across the mixed stream to $5.93$\,ms, delivering an effective throughput of $168.6$\,FPS and a measured $1.91\times$ end-to-end acceleration.
\end{enumerate}
This confirms that the dual-branch peak-pooling triage head provides genuine clinical screening efficiency with rigorous mathematical safety guarantees, directly solving the open limitation of prior all-positive benchmarks.

\subsection{Qualitative Analysis}
Fig.~\ref{fig:qualitative} compares representative predictions. U-Net and UNet++ exhibit boundary leakage into mucosal folds and are distracted by specular fluid reflections; Attention U-Net produces over-smoothed contours that miss diminutive protrusions. FEE-Net reproduces sharp, anatomically faithful boundaries while suppressing glare-induced false edges, consistent with the quantitative HD95 reduction.

\begin{figure*}[t]
\centering
\IfFileExists{figure_qualitative_comparison.png}{%
\includegraphics[width=0.95\textwidth]{figure_qualitative_comparison.png}%
}{%
\fbox{\parbox[c][4.5cm][c]{0.92\textwidth}{\centering\footnotesize
\texttt{figure\_qualitative\_comparison.png} not present in this build.\\ It is regenerated by the benchmark notebook and shipped in the results archive.}}%
}
\caption{\textbf{Qualitative comparison on independent test cases.} From left to right: endoscopic frame, expert ground truth, U-Net, Attention U-Net, and the proposed FEE-Net. FEE-Net adheres closely to lesion boundaries and resists specular glare artifacts.}
\label{fig:qualitative}
\end{figure*}

\subsection{Limitations}
We identify several limitations. (1) \emph{Baselines:} our comparison covers the classical clinical baselines (U-Net, Attention U-Net, UNet++) under identical budgets; specialized polyp networks (e.g., PraNet~\cite{fan2020pranet}) and transformer/foundation models were not retrained here, so no claim of superiority over them is made. (2) \emph{Statistical scope and protocol evolution:} each model was evaluated on the isolated 150-case independent test split with a fixed seed. While an earlier v2.1 prototype under a legacy 5-fold cross-validation protocol recorded $81.69\pm1.28\%$ Dice, that evaluation lacked strict patient-level partition guarantees and utilized a simplified inline backbone; under the rigorous, zero-leakage protocol with the fully aligned repository architecture, FEE-Net yields $79.47\%$ (and $81.42\%$ with early-exit triage multi-task regularization), which we deliberately report to uphold scientific honesty against optimistic leakage artifacts. (3) \emph{Mucosa crop approximation:} normal mucosal samples in the mixed cohort are derived from non-lesion regions of polyp frames; full-procedure screening streams with varying lighting and debris present further distribution shifts. (4) \emph{Hardware profiling scope:} latency measurements were conducted on a dedicated NVIDIA RTX 4060 Laptop GPU, which reflects edge clinic workstations but differs from multi-GPU cloud environments. (5) \emph{Single dataset:} all experiments use Kvasir-SEG; cross-dataset generalization (e.g., CVC-ClinicDB, ETIS) is untested. (6) \emph{Component-level novelty:} the wavelet decomposition, channel attention, depthwise directional filtering, bilinear upsampling, deep supervision, and the boundary-ribbon loss are established components (Sec.~II); our contribution is their specific integration and the certified screening protocol, not a novel component in isolation. (7) \emph{Cohort scale and bounds:} while the current mixed cohort achieves an exact finite-sample upper bound of $3.92\%$ at $\tau^*=0.59$, compressing the bound to stringent clinical safety thresholds ($\le 1.0\%$) requires scaling the calibration and certification cohorts to multi-center video repositories.

\section{Conclusion}
This paper introduced an integrated, safety-certified framework for colorectal polyp detection and screening, combining the lightweight dual-domain FEE-Net architecture (8.539M parameters) with finite-sample Conformal Risk Control. By coupling 2D Haar wavelet decomposition with Spectral Channel Attention and Omni-Directional Wavelet Filtering inside an edge-gated decoder, FEE-Net targets optical specular glare and directional boundary reconstruction while maintaining high computational efficiency; a dual-branch local-peak triage head with stop-gradient probing screens non-pathological frames without degrading dense segmentation representations. Calibrated via the exact Clopper--Pearson binomial bound under strict separation of calibration and certification, the triage mechanism was evaluated on a leakage-free 417-sample mixed cohort derived from the held-out test split, certifying $\mathrm{FNR} \le 3.92\% < 5.0\%$ with $95\%$ confidence, achieving a measured $0.00\%$ false-negative rate, and safely bypassing $99.25\%$ of healthy mucosal frames (providing an effective $1.91\times$ end-to-end inference acceleration, reducing latency from $11.30$\,ms to $5.93$\,ms). Rigorous evaluations under a strict zero-leakage protocol establish that clinical utility in high-throughput endoscopy is governed by statistical safety guarantees and transparent reproducible protocols rather than unconstrained cross-literature SOTA claims.

\bibliographystyle{IEEEtran}
\begin{thebibliography}{00}

\bibitem{ronneberger2015u} O. Ronneberger, P. Fischer, and T. Brox, ``U-Net: Convolutional networks for biomedical image segmentation,'' in \emph{Proc. MICCAI}, 2015, pp. 234--241.

\bibitem{oktay2018attention} O. Oktay \emph{et al.}, ``Attention U-Net: Learning where to look for the pancreas,'' in \emph{Proc. MIDL}, 2018.

\bibitem{zhou2019unetplusplus} Z. Zhou, M. M. Rahman Siddiquee, N. Tajbakhsh, and J. Liang, ``UNet++: A nested U-Net architecture for medical image segmentation,'' \emph{IEEE Trans. Med. Imag.}, vol. 39, no. 1, pp. 3--11, 2020.

\bibitem{chen2021transunet} J. Chen \emph{et al.}, ``TransUNet: Transformers make strong encoders for medical image segmentation,'' \emph{arXiv preprint arXiv:2102.04306}, 2021.

\bibitem{ma2024medsam} J. Ma \emph{et al.}, ``Segment anything in medical images,'' \emph{Nature Communications}, vol. 15, no. 1, art. 654, 2024.

\bibitem{ruan2024vmunet} J. Ruan and S. Xiang, ``VM-UNet: Vision Mamba UNet for medical image segmentation,'' \emph{arXiv preprint arXiv:2402.02491}, 2024.

\bibitem{fan2020pranet} D.-P. Fan, G.-P. Ji, T. Zhou, G. Chen, H. Fu, J. Shen, and L. Shao, ``PraNet: Parallel reverse attention network for polyp segmentation,'' in \emph{Proc. MICCAI}, 2020, pp. 263--273.

\bibitem{hu2018senet} J. Hu, L. Shen, and G. Sun, ``Squeeze-and-excitation networks,'' in \emph{Proc. IEEE/CVF CVPR}, 2018, pp. 7132--7141.

\bibitem{odena2016deconv} A. Odena, V. Dumoulin, and C. Olah, ``Deconvolution and checkerboard artifacts,'' \emph{Distill}, vol. 1, no. 10, 2016.

\bibitem{bates2021distribution} S. Bates, A. N. Angelopoulos, L. Lei, J. Malik, and M. I. Jordan, ``Distribution-free, risk-controlling prediction sets,'' \emph{Journal of the ACM}, vol. 68, no. 6, pp. 1--34, 2021.

\bibitem{angelopoulos2022crc} A. N. Angelopoulos, S. Bates, A. Fisch, L. Lei, and T. Schuster, ``Conformal risk control,'' \emph{arXiv preprint arXiv:2208.02814}, 2022.

\bibitem{clopper1934use} C. J. Clopper and E. S. Pearson, ``The use of confidence or fiducial limits illustrated in the case of the binomial,'' \emph{Biometrika}, vol. 26, no. 4, pp. 404--413, 1934.

\bibitem{angelopoulos2021gentle} A. N. Angelopoulos and S. Bates, ``A gentle introduction to conformal prediction and distribution-free uncertainty quantification,'' \emph{arXiv preprint arXiv:2107.07511}, 2021.

\bibitem{eljurdi2021perimeter} R. El Jurdi, C. Petitjean, P. Honeine, V. Cheplygina, and F. Abdallah, ``A surprisingly effective perimeter-based loss for medical image segmentation,'' in \emph{Proc. MIDL}, 2021, pp. 158--167.

\bibitem{karimi2019hd} D. Karimi and S. E. Salcudean, ``Reducing the Hausdorff distance in medical image segmentation with convolutional neural networks,'' \emph{IEEE Trans. Med. Imag.}, vol. 39, no. 2, pp. 365--376, 2020.

\bibitem{karimi2022boundaryiou} D. Karimi and S. E. Salcudean, ``Boundary IoU: Improving object-centric image segmentation evaluation,'' \emph{Medical Image Analysis}, vol. 78, art. 102301, 2022.

\bibitem{zhao2023wranet} Y. Zhao, S. Wang, Y. Zhang, S. Qiao, and M. Zhang, ``WRANet: Wavelet integrated residual attention U-Net network for medical image segmentation,'' \emph{Complex \& Intelligent Systems}, vol. 9, no. 6, pp. 6971--6983, 2023.

\bibitem{rajasekar2024adaptunet} D. Rajasekar, G. Theja, M. R. Prusty, and S. Chinara, ``Efficient colorectal polyp segmentation using wavelet transformation and AdaptUNet: A hybrid U-Net,'' \emph{Heliyon}, vol. 10, no. 13, art. e33655, 2024.

\bibitem{dai2025crcfdr} M. Dai, W. Luo, and T. Li, ``Statistical management of the false discovery rate in medical instance segmentation based on conformal risk control,'' in \emph{Proc. IEEE ICIPCA}, 2025.

\bibitem{cls2024} B. Tan, Z. Wang, J. Duan, K. Xu, H. T. Shen, F. Shen, and X. Shi, ``Conformal 3D lesion segmentation with balanced risk control,'' in \emph{Proc. MICCAI}, 2024.

\bibitem{zheng2026crc} D. Zheng \emph{et al.}, ``What can conformal risk control certify for PET/CT tumour segmentation?'' in \emph{Proc. AISTATS}, 2026.

\bibitem{jha2020kvasir} D. Jha, P. H. Smedsrud, M. A. Riegler, P. Halvorsen, T. de Lange, D. Johansen, and H. D. Johansen, ``Kvasir-SEG: A segmented polyp dataset,'' in \emph{Proc. Int. Conf. Multimedia Model. (MMM)}, 2020, pp. 451--462.

\end{thebibliography}

\end{document}
